\chapter{Integral Geometry}

\begin{center}\itshape
How many lines meet a curve? The question needs a way to measure sets of lines, and symmetry allows exactly one. With it, Buffon's needle becomes
Crofton's formula, the length of any curve becomes the average number of times a random line crosses it, and Bertrand's famous paradox finds its answer.
\end{center}

In 1889 Joseph Bertrand, permanent secretary of the Paris Academy of Sciences, opened his \emph{Calcul des probabilités} with a warning. Draw a chord at random
in a circle: what is the probability that it is longer than the side of the inscribed equilateral triangle? Choose its two endpoints at random on the circle,
and the answer is $\frac13$; choose its distance to the centre at random along a radius, and it is $\frac12$; choose its midpoint at random in the disc, and it is
$\frac14$. Volume~4 computed all three. Bertrand's conclusion was that the question is ill posed: ``none of the three is false, none is correct'', because ``at random''
means nothing until one says how.

Henri Poincaré took up the question in his own \emph{Calcul des probabilités} of 1896. Among the ways of choosing a line, he observed, one is distinguished: the measure
$\dd p\,\dd\theta$, in the coordinates where $p$ is the distance of the line to an origin and $\theta$ the direction of its normal, is the only one that does not change
when the whole figure is moved. Seventy-seven years later the physicist Edwin Jaynes, a champion of probability as logic, turned this into an argument and an experiment.
If straws are tossed onto a circle drawn on the floor by someone who does not aim, the answer cannot depend on exactly where the circle is drawn, or how it is turned;
only one distribution of lines has that property, and it gives $\frac12$. Jaynes reported that tossed straws agreed. This chapter proves the uniqueness that
Poincaré noticed, and draws its consequences: Crofton's formula for the length of a curve, Cauchy's formulas for perimeters and surface areas, and Poincaré's own
formula for the intersections of two moving curves.

The subject had begun with Buffon's needle (1777), and Volume~9, Appendix~J proved the Buffon--Barbier theorem for polygonal ``noodles'', Cauchy's formula for the
perimeter of a convex set, and Barbier's theorem on curves of constant width. Here these become special cases of a single principle, valid for every piecewise $C^1$ curve.

\section{The measure of lines}

An \emph{oriented line} in the plane is determined by its unit normal $u_\theta=(\cos\theta,\sin\theta)$, $\theta\in\R/2\pi\Z$, and the number $p$ with
\[\ell(\theta,p)=\{x\in\R^2:\langle x,u_\theta\rangle=p\},\]
oriented by $u_\theta$ turned by $+\frac\pi2$. This identifies the set $\mathcal L^+$ of oriented lines with the cylinder $(\R/2\pi\Z)\times\R$. Forgetting the orientation identifies
$(\theta,p)$ with $(\theta+\pi,-p)$, and the set $\mathcal L$ of lines is parametrized once by $\theta\in[0,\pi)$, $p\in\R$. A rigid motion acts on lines: the rotation by $\alpha$ about the
origin maps $\ell(\theta,p)$ to $\ell(\theta+\alpha,p)$, and the translation by $v$ maps it to $\ell\bigl(\theta,p+\langle v,u_\theta\rangle\bigr)$. Since rotations and translations generate the
motions, the measure $\dd\theta\,\dd p$ is invariant under all of them: rotations shift $\theta$, and for each fixed $\theta$ a translation shifts $p$ (Fubini).

\begin{theorem}[uniqueness of the measure on lines]\label{thm:linesunique}
Let $\nu$ be a Borel measure on $\mathcal L^+$, finite on compact sets and invariant under rotations and translations. Then $\nu=c\,\dd\theta\,\dd p$ for a constant $c\ge0$. Likewise every
motion-invariant Borel measure on $\mathcal L$, finite on compact sets, is $c\,\dd\theta\,\dd p$ on $[0,\pi)\times\R$.
\end{theorem}

\begin{proof}
\emph{Step 1: invariance under uniform shifts of $p$.} Let $a$ be bounded Borel on the circle, $b$ continuous with compact support on $\R$, and $t\in\R$; we show that
$\nu\bigl(a(\theta)b(p+t)\bigr)=\nu\bigl(a(\theta)b(p)\bigr)$, writing $\nu(f)$ for $\int f\dd\nu$. Cut the circle into arcs $I_1,\dots,I_m$ of length $\delta=\frac{2\pi}m$, with centres $\theta_j$. The translation by
$tu_{\theta_j}$ maps $\ell(\theta,p)$ to $\ell(\theta,p+t\cos(\theta-\theta_j))$, so by invariance
\[\nu\bigl(a\mathbf 1_{I_j}(\theta)\,b(p+t\cos(\theta-\theta_j))\bigr)=\nu\bigl(a\mathbf 1_{I_j}(\theta)\,b(p)\bigr).\]
For $\theta\in I_j$, $|t\cos(\theta-\theta_j)-t|\le|t|\frac{\delta^2}8$, so $|b(p+t\cos(\theta-\theta_j))-b(p+t)|\le\omega(|t|\delta^2)$, where $\omega$ is the modulus of continuity of $b$. All the functions involved vanish
outside $(\R/2\pi\Z)\times[-R,R]$ for some $R$ independent of $\delta$. Summing over $j$,
\[\bigl|\nu\bigl(a(\theta)b(p+t)\bigr)-\nu\bigl(a(\theta)b(p)\bigr)\bigr|\le\|a\|_\infty\,\omega(|t|\delta^2)\,\nu\bigl((\R/2\pi\Z)\times[-R,R]\bigr),\]
and letting $\delta\to0$ gives the claim.

\emph{Step 2: the product structure.} A Borel measure $m$ on $\R$, finite on compact sets and invariant under all translations, is $\kappa$ times Lebesgue measure with $\kappa=m([0,1))$: the $n$ translates
$[\frac kn,\frac{k+1}n)$ have equal measure, so $m([0,\frac1n))=\frac\kappa n$, hence $m(J)=\kappa|J|$ for intervals with rational endpoints, then for all intervals by monotone limits, and these form a $\pi$-system
generating the Borel sets (Volume~7, Theorem~7.5). The same argument applies to rotation-invariant finite measures on the circle, with arcs of length $\frac{2\pi}n$.

Fix an arc $I$ and let $m_I(B)=\nu(I\times B)$ for Borel $B\subset\R$. By Step~1 with $a=\mathbf 1_I$, $\int b(p+t)\dd m_I=\int b\dd m_I$ for continuous $b$ with compact support, and approximating indicators of compact intervals by
such functions (dominated convergence), $m_I$ is translation invariant. So $\nu(I\times J)=\kappa(I)|J|$ for intervals $J$. Now $I\mapsto\kappa(I)=\nu(I\times[0,1))$ extends to a finite measure on the circle, invariant under
rotations because $\nu$ is, so $\kappa(I)=c|I|$. Hence $\nu(I\times J)=c|I||J|$, and these products form a $\pi$-system generating the Borel sets of the cylinder: $\nu=c\,\dd\theta\,\dd p$.

\emph{Unoriented lines.} Given $\nu$ on $\mathcal L$, let $\tilde\nu(f)=\nu(\bar f)$ on $\mathcal L^+$, where $\bar f(\ell)$ is the sum of $f$ over the two orientations of $\ell$. Motions commute with forgetting the
orientation, so $\tilde\nu$ is invariant, hence $\tilde\nu=c\,\dd\theta\,\dd p$. For $F$ on $\mathcal L$, apply this to $f=F\circ\pi$, with $\pi$ the forgetful map: $2\nu(F)=c\int_0^{2\pi}\!\!\int F\circ\pi\dd p\dd\theta=2c\int_0^\pi\!\!\int F\dd p\dd\theta$.
\end{proof}

From now on, $\dd\ell=\dd p\,\dd\theta$ on $[0,\pi)\times\R$ is \emph{the} measure on lines. It is the image of the Haar measure of the group of motions (Chapter~8) under $g\mapsto g(\ell_0)$, but we
did not need that.

\begin{corollary}[Bertrand's chord, answered]\label{cor:bertrand}
The lines meeting the unit disc have measure $2\pi$. For a line chosen according to the normalized measure among them, the distance $|p|$ to the centre is uniform on $[0,1]$, and the chord
is longer than $\sqrt3$ with probability $\frac12$.
\end{corollary}

\begin{proof}
The line $\ell(\theta,p)$ meets the disc iff $|p|\le1$: measure $\pi\cdot2$. Under the normalized measure, $\theta$ and $p$ are independent and uniform, so $|p|$ is uniform on $[0,1]$. The chord has length
$2\sqrt{1-p^2}>\sqrt3$ iff $|p|<\frac12$.
\end{proof}

The two other answers of Bertrand correspond to laws on chords that are not restrictions of an invariant measure on lines: with random endpoints, $|p|$ has density $\frac2\pi(1-p^2)^{-1/2}$,
and with a random midpoint, density $2p$ (Exercise~\ref{ex:bertrandlaws}). Jaynes's point was that ``a random line falling on the circle'', thrown without aiming at the circle, must be the
restriction of a motion-invariant law on all lines; Theorem~\ref{thm:linesunique} leaves no choice.

\section{Crofton's formula}

For a curve $c$ and a line $\ell$, let $n(c\cap\ell)$ be the number of parameters $t$ with $c(t)\in\ell$ (points where the curve crosses itself on $\ell$ count with multiplicity).

\begin{lemma}[counting preimages]\label{lem:banach}
Let $g\colon[0,1]\to\R$ be $C^1$. Then $p\mapsto\#g^{-1}(p)$ is a Borel function and
\[\int_\R\#g^{-1}(p)\dd p=\int_0^1|g'(t)|\dd t.\]
\end{lemma}

\begin{proof}
The open set $\{t\in(0,1):g'(t)\neq0\}$ is a countable disjoint union of open intervals $J_k$. On each, $g$ is strictly monotone, so $\#(g^{-1}(p)\cap J_k)$ is the indicator of the interval $g(J_k)$, whose length is
$\int_{J_k}|g'|$ by the fundamental theorem of calculus. Every other parameter lies in $\{0,1\}$ or in the critical set $Z=\{g'=0\}$, and $g(Z)$ is null by Sard's theorem in dimension $1$
(Theorem~\ref{thm:sard}). So for almost every $p$, $\#g^{-1}(p)=\sum_k\mathbf 1_{g(J_k)}(p)$, and by Tonelli $\int\#g^{-1}=\sum_k\int_{J_k}|g'|=\int_0^1|g'|$, since $g'=0$ on $Z$. (The function $\#g^{-1}$ is the increasing
limit, as $m\to\infty$, of the number of intervals $[\frac j{2^m},\frac{j+1}{2^m})$, together with $\{1\}$, whose image contains $p$, each the indicator of an interval; so it is Borel.)
\end{proof}

\begin{theorem}[Crofton, 1868]\label{thm:crofton}
For a piecewise $C^1$ curve $c$ in the plane of length $L$,
\[\int_{\mathcal L}n(c\cap\ell)\dd\ell=2L.\]
In words: the length of a curve is half the measure of the lines that meet it, counted with multiplicity.
\end{theorem}

\begin{proof}
Both sides are additive over the $C^1$ pieces (points where pieces join lie on a null set of lines), so let $c\colon[0,1]\to\R^2$ be $C^1$. For fixed $\theta$, $n(c\cap\ell(\theta,p))$ is the number of preimages of $p$ by
$g_\theta(t)=\langle c(t),u_\theta\rangle$, so by Lemma~\ref{lem:banach} and Tonelli,
\[\int_{\mathcal L}n\dd\ell=\int_0^\pi\!\!\int_0^1|\langle c'(t),u_\theta\rangle|\dd t\dd\theta=\int_0^1|c'(t)|\int_0^\pi|\cos(\theta-\varphi(t))|\dd\theta\dd t,\]
where $\varphi(t)$ is the direction of $c'(t)$ (where $c'(t)\neq0$). The inner integral is $2$ for every $\varphi$, and $\int_0^1|c'|=L$.
\end{proof}

\begin{figure}[ht]
\centering
\fig{ch10-crofton}
\caption{Left: random lines for the invariant measure, among those meeting a disc, and a curve inside it; the crossings are marked. Averaged over many lines, the number of crossings is
$\frac{2L}{2\pi R}$ by Crofton's formula, where $2\pi R$ is the measure of the lines meeting the disc. Right: a Reuleaux triangle between pairs of parallel support lines, at the same distance $w$ in every direction.}
\label{fig:crofton}
\end{figure}

\begin{corollary}[Buffon's noodle]\label{cor:noodle}
A piecewise $C^1$ curve of length $L$, thrown on parallel lines at distance $d$ with uniform random orientation and uniform random offset, meets them on average $\frac{2L}{\pi d}$ times.
\end{corollary}

\begin{proof}
Fix the curve and move the lines instead: their common normal direction $\theta$ is uniform on $[0,\pi)$ and the offset $p_0$ uniform on $[0,d)$, the lines being $\ell(\theta,p_0+kd)$, $k\in\Z$. The expected number of crossings is
\[\frac1{\pi d}\int_0^\pi\!\!\int_0^d\sum_kn\bigl(c\cap\ell(\theta,p_0+kd)\bigr)\dd p_0\dd\theta=\frac1{\pi d}\int_{\mathcal L}n\dd\ell=\frac{2L}{\pi d}.\]
\end{proof}

For a needle of length $L\le d$, which meets at most one line, this is the probability of crossing: Buffon's $\frac{2L}{\pi d}$ of Volume~4, Chapter~1, now for any curve. Crofton's formula is
also a measuring instrument: to estimate the length of a curve in an image, count its crossings with a family of lines in a few directions and multiply. Hugo Steinhaus built such a
transparent ``longimeter'', and biologists measure the length of roots and blood vessels on microscope slides in the same way.

\section{Convex sets}

Let $K$ be a compact convex set in the plane with non-empty interior, whose boundary is a piecewise $C^1$ closed curve of length $L(K)$, and let $w_K(\theta)$ be its width in the direction $u_\theta$, the
length of the interval $\{\langle x,u_\theta\rangle:x\in K\}$.

\begin{lemma}\label{lem:convexline}
If $x_0\in\operatorname{int}K$ and $y\in K$, the segment $[x_0,y)$ lies in $\operatorname{int}K$. Consequently a line through an interior point meets $\partial K$ in exactly two points, and the lines meeting $K$ but not
its interior form a null set.
\end{lemma}

\begin{proof}
If the ball $B(x_0,r)$ lies in $K$ and $z=(1-\lambda)x_0+\lambda y$ with $0\le\lambda<1$, then $B(z,(1-\lambda)r)=(1-\lambda)B(x_0,r)+\lambda y\subset K$ by convexity. A line through $x_0$ meets $K$ in a segment $[a,b]$ containing $x_0$; by
the first part, the points of $(a,b)$ are interior, while $a$ and $b$, with points outside $K$ arbitrarily close on the line, are on the boundary. For each $\theta$, the lines $\ell(\theta,p)$ meeting $K$ have $p$ in a compact interval
$[h_-(\theta),h_+(\theta)]$ of length $w_K(\theta)$, and the projection of the open convex set $\operatorname{int}K$ is an open interval dense in it (since $K$ is the closure of its interior, by the first part); so
only the two values $p=h_\pm(\theta)$ give lines missing the interior.
\end{proof}

\begin{theorem}[Crofton and Cauchy]\label{thm:cauchycrofton}
The lines meeting $K$ have measure $L(K)$, and
\[L(K)=\int_0^\pi w_K(\theta)\dd\theta.\]
\end{theorem}

\begin{proof}
By Lemma~\ref{lem:convexline}, $n(\partial K\cap\ell)$ equals $2$ for almost every line meeting $K$, and $0$ for the others. Crofton's formula gives $2L(K)=2\int\mathbf 1_{\ell\cap K\neq\emptyset}\dd\ell=2\int_0^\pi w_K(\theta)\dd\theta$.
\end{proof}

This is Cauchy's formula of Volume~9 (Theorem~11.5), proved there by inscribed polygons; Barbier's theorem (a set of constant width $w$ has perimeter $\pi w$) follows in one line, as in Volume~9,
Corollary~K.6. The new information is the first statement: \emph{the measure of the lines meeting a convex set is its perimeter}. Ratios of perimeters are therefore probabilities. If $K'\subset K$ are
convex, a random line meeting $K$ meets $K'$ with probability $\frac{L(K')}{L(K)}$, whatever the shapes. Exercises~\ref{ex:chord} and~\ref{ex:twolines} continue in this vein; the second shows that two random
lines meeting a convex set cross inside it with probability $\frac{2\pi A}{L^2}$, at most $\frac12$ by the isoperimetric inequality of Volume~9, with equality only for the disc.

\section{Moving curves and shadows}

Crofton's formula moves lines. Poincaré's formula moves a whole curve. A rigid motion $g$ is the rotation by $\alpha\in[0,2\pi)$ followed by the translation by $v\in\R^2$; on the group of motions we use the
measure $\dd g=\dd v\,\dd\alpha$, its Haar measure (it is invariant on both sides: left multiplication by a translation shifts $v$, by a rotation rotates $v$ and shifts $\alpha$, both preserving $\dd v\,\dd\alpha$, and
similarly on the right).

\begin{theorem}[Poincaré, 1896]\label{thm:poincare}
Let $C_1$, $C_2$ be polygonal curves of lengths $L_1$, $L_2$. Then
\[\int n(C_1\cap gC_2)\dd g=4L_1L_2.\]
\end{theorem}

\begin{proof}
For almost every $g$, no vertex of one curve lies on the other, and $n(C_1\cap gC_2)$ is the sum over pairs of sides $(s,s')$ of $n(s\cap gs')\in\{0,1\}$ (two segments meet at most once unless they overlap, which happens for
$g$ in a null set). So it suffices to treat two segments, of lengths $a$ and $b$. For a fixed rotation, $s\cap(R_\alpha s'+v)\neq\emptyset$ iff $v\in s-R_\alpha s'$, a parallelogram with sides $a$ and $b$, of area $ab|\sin\beta|$,
where $\beta$ is the angle between $s$ and $R_\alpha s'$, which is $\alpha$ plus a constant. So $\int n(s\cap gs')\dd g=\int_0^{2\pi}ab|\sin\beta|\dd\alpha=4ab$, and summing over pairs gives $4L_1L_2$.
\end{proof}

The formula extends to rectifiable curves by approximation. The full \emph{kinematic formula} of Blaschke, Santaló and Chern measures the motions
$g$ for which two convex sets $K_1$ and $gK_2$ meet: $2\pi(A_1+A_2)+L_1L_2$.

In space, the role of the width is played by the area of the shadow. For a unit vector $u$ let $\pi_u$ be the orthogonal projection onto the plane $u^\perp$, and let $\sigma$ be the uniform probability on $S^2$, the normalized
area measure. \emph{Archimedes' hat-box theorem} says that for $u$ distributed according to $\sigma$, the coordinate $\langle u,n\rangle$ along any unit vector $n$ is uniform on $[-1,1]$: with $n$ the north pole and
$u=(\sqrt{1-z^2}\cos\phi,\sqrt{1-z^2}\sin\phi,z)$, the partial derivatives $\partial_\phi u$ and $\partial_zu$ are orthogonal with lengths $\sqrt{1-z^2}$ and $\frac1{\sqrt{1-z^2}}$, so the area element is $\dd z\,\dd\phi$. In particular
$\mathbb E|\langle u,n\rangle|=\frac12$.

\begin{theorem}[Cauchy's surface area formula]\label{thm:cauchyspace}
Let $K\subset\R^3$ be a convex polytope with non-empty interior and surface area $A$. Then the mean area of its shadows is a quarter of its surface area:
\[\int_{S^2}\operatorname{area}\bigl(\pi_u(K)\bigr)\dd\sigma(u)=\frac A4.\]
\end{theorem}

\begin{proof}
Let the faces $F$ have areas $A_F$ and outer unit normals $n_F$. On the plane of $F$, $\pi_u$ is an affine map multiplying areas by $|\langle n_F,u\rangle|$. Fix $u$ with $\langle n_F,u\rangle\neq0$ for every face $F$ (this excludes finitely many great circles of directions, a null set). The projections of the edges are segments, a null set. If $y\in\pi_u(K)$ is not on them and is interior to $\pi_u(K)$, the line $y+\R u$ meets $\operatorname{int}K$ (the projection of $\operatorname{int}K$ is an
open convex set whose closure is $\pi_u(K)$, and an open convex set in the plane contains the interior of its closure, since a point outside it can be separated from it by a line); by the three-dimensional version of
Lemma~\ref{lem:convexline}, which has the same proof, this line meets $\partial K$ in exactly two points, both interior to faces. The boundary of $\pi_u(K)$ is null too. Therefore
\[\sum_FA_F|\langle n_F,u\rangle|=\sum_F\operatorname{area}\pi_u(F)=\int_{u^\perp}\#\bigl(\pi_u^{-1}(y)\cap\partial K\bigr)\dd y=2\operatorname{area}\pi_u(K).\]
Averaging over $u$ with $\mathbb E|\langle n_F,u\rangle|=\frac12$ gives $\frac12\sum_FA_F\cdot\frac12=\frac A4$.
\end{proof}

\begin{beauty}
One single number, the measure of lines or of planes, measures length, perimeter and area, and it does so because it is the only measure that the symmetries of space allow. Hugo Hadwiger
proved in 1957 the definitive form of this idea: every function on convex bodies in $\R^n$ that is additive (in the sense $\mu(K\cup K')+\mu(K\cap K')=\mu(K)+\mu(K')$), continuous and invariant under
motions is a linear combination of $n+1$ quantities, the \emph{intrinsic volumes}: in the plane, the area, the perimeter and the Euler characteristic. Volume, surface area, mean width and
Euler's $\chi$ of Chapter~6 are the only measurements that symmetry permits. Integral geometry works just as well in the hyperbolic plane of Chapter~9, with the invariant measure on its geodesics,
and its descendants, the Radon transform and its inversion, are the mathematics of the medical scanner.
\end{beauty}

\section{Exercises}

\begin{exercise}[$\star$]\label{ex:bertrandlaws}
For Bertrand's methods (a) random endpoints and (c) random midpoint, compute the law of the distance $|p|$ of the chord to the centre, and show that neither law makes $(\theta,p)$ uniform. Under the
invariant measure, compute the law and the mean of the chord length.
\end{exercise}

\begin{exercise}[$\star$]\label{ex:chord}
\emph{Crofton's chord formula.} For a convex $K$ of area $A$, let $\sigma(\ell)$ be the length of $\ell\cap K$. Show that $\int\sigma(\ell)\dd\ell=\pi A$. Deduce that the mean chord of a random line meeting $K$ is
$\frac{\pi A}{L}$, and compute it for a disc and a square.
\end{exercise}

\begin{exercise}[$\star\star$]\label{ex:twolines}
Two independent random lines meet a convex set $K$ (each with the normalized invariant measure among the lines meeting $K$). Show that they cross inside $K$ with probability $\frac{2\pi A}{L^2}$. Evaluate
it for a disc and a square, and show that it is at most $\frac12$.
\end{exercise}

\begin{exercise}[$\star\star$]\label{ex:laplacegrid}
\emph{Laplace's grid} (1812). A needle of length $l\le d$ is thrown on a floor of square tiles of side $d$. Show that it crosses a joint with probability $\frac{4ld-l^2}{\pi d^2}$.
\end{exercise}

\begin{exercise}[$\star$]\label{ex:shadows}
(a) Show that the mean shadow area of a unit cube is $\frac32$, and that its largest shadow, seen along a diagonal, is a regular hexagon of area $\sqrt3$. (b) What is the mean shadow of a sphere
of radius $r$, and is it consistent with Theorem~\ref{thm:cauchyspace}?
\end{exercise}

\begin{exercise}[$\star\star$]\label{ex:spherecrofton}
\emph{Crofton on the sphere.} Each great circle of $S^2$ is $u^\perp\cap S^2$ for a pole $u$, determined up to sign. For a $C^1$ curve on $S^2$ of length $L$, show that
$\int_{S^2}n(c\cap u^\perp)\dd\sigma(u)=\frac L\pi$. Check it on a great circle. (This is how Steinhaus proposed to measure the length of a curve on a globe.)
\end{exercise}

\section{Hints}

\begin{hint}{ex:bertrandlaws}
For (a), with one endpoint at angle $0$ and the other at angle $\varphi$, $|p|=|\cos\frac\varphi2|$.
\end{hint}

\begin{hint}{ex:chord}
For fixed $\theta$, integrate the chord length over $p$: Cavalieri.
\end{hint}

\begin{hint}{ex:twolines}
Fix the first line; the second crosses its chord with a measure given by Theorem~\ref{thm:cauchycrofton} applied to a segment. Then use Exercise~\ref{ex:chord}.
\end{hint}

\begin{hint}{ex:laplacegrid}
Count horizontal and vertical crossings separately, and compute the probability of crossing both.
\end{hint}

\begin{hint}{ex:shadows}
A face of the cube with normal $n$ contributes $|\langle n,u\rangle|$ to the sum in the proof.
\end{hint}

\begin{hint}{ex:spherecrofton}
For an arc element at $x$ with unit tangent $t$, the great circle $u^\perp$ crosses it when $u$ lies in a thin band around the great circle $x^\perp$, of angular width $|\langle u,t\rangle|\dd s$.
\end{hint}

\section{Solutions}

\begin{solution}{ex:bertrandlaws}
\idea Each method induces a law on $|p|$; the invariant measure is the one with $|p|$ uniform.
\writeup (a) With $\varphi$ uniform on $[0,2\pi)$, $|p|=|\cos\frac\varphi2|$, and $\frac\varphi2$ is uniform on $[0,\pi)$; so $|p|=|\cos\psi|$ with $\psi$ uniform, which has density $\frac2\pi(1-p^2)^{-1/2}$ on $[0,1]$: chords far from the centre are
favoured. (c) The midpoint is uniform in the disc and $|p|$ is its distance to the centre: $\mathbb P(|p|\le r)=r^2$, density $2p$: chords near the edge again, more strongly. Under the invariant measure $|p|$ is uniform, so the
chord length $\sigma=2\sqrt{1-p^2}$ has $\mathbb P(\sigma>s)=\mathbb P(|p|<\sqrt{1-s^2/4})=\sqrt{1-s^2/4}$ for $0\le s\le2$, and $\mathbb E\sigma=\int_0^12\sqrt{1-p^2}\dd p=\frac\pi2$.
\end{solution}

\begin{solution}{ex:chord}
\idea The chords in one direction sweep out the area.
\writeup For fixed $\theta$, $\int\sigma(\ell(\theta,p))\dd p$ is the area of $K$, by Fubini in coordinates along $u_\theta$ and perpendicular to it. Integrating over $\theta\in[0,\pi)$ gives $\pi A$. The lines meeting $K$ have measure
$L$, so the mean chord is $\frac{\pi A}L$. For the disc of radius $r$: $\frac{\pi\cdot\pi r^2}{2\pi r}=\frac{\pi r}2$, in agreement with Exercise~\ref{ex:bertrandlaws}. For the unit square: $\frac\pi4\approx0.785$, less than the side, because many
lines cut off a corner.
\end{solution}

\begin{solution}{ex:twolines}
\idea The measure of the lines crossing a segment of length $\sigma$ is $2\sigma$.
\writeup By Theorem~\ref{thm:crofton} applied to a segment of length $\sigma$ (a line meets it at most once, except for a null set), the lines meeting it have measure $2\sigma$. Two lines meeting $K$ cross inside $K$ iff the second
meets the chord $\ell_1\cap K$. So the measure of such pairs is $\int2\sigma(\ell_1)\dd\ell_1=2\pi A$ (Exercise~\ref{ex:chord}), while all pairs of lines meeting $K$ have measure $L^2$. The probability is $\frac{2\pi A}{L^2}$. For the disc,
$\frac{2\pi\cdot\pi r^2}{4\pi^2r^2}=\frac12$; for the square, $\frac{2\pi}{16}\approx0.39$. The isoperimetric inequality $4\pi A\le L^2$ (Volume~9, Corollary~J.8) gives the bound $\frac12$, with equality only for discs.
\end{solution}

\begin{solution}{ex:laplacegrid}
\idea Inclusion--exclusion: horizontal crossings plus vertical crossings minus both.
\writeup Let $N_h$, $N_v$ be the numbers of crossings with the horizontal and vertical lines; each is $0$ or $1$ (almost surely, since $l\le d$), and by Corollary~\ref{cor:noodle} each has mean $\frac{2l}{\pi d}$. Given the angle $\theta$ of
the needle with the horizontal, its horizontal and vertical extents are $l|\cos\theta|$ and $l|\sin\theta|$, and the two offsets of the grid are independent and uniform, so by the argument of Volume~9, Lemma~J.3,
$\mathbb P(N_h=N_v=1\mid\theta)=\frac{l^2|\sin\theta\cos\theta|}{d^2}$. Averaging over $\theta$ uniform on $[0,2\pi)$, $\mathbb P(N_h=N_v=1)=\frac{l^2}{d^2}\cdot\frac1{2\pi}\int_0^{2\pi}\frac{|\sin2\theta|}2\dd\theta=\frac{l^2}{\pi d^2}$. Hence
$\mathbb P(N_h+N_v\ge1)=\frac{4l}{\pi d}-\frac{l^2}{\pi d^2}=\frac{4ld-l^2}{\pi d^2}$. For $l=d$ it is $\frac3\pi\approx0.95$.
\end{solution}

\begin{solution}{ex:shadows}
\idea The cube has surface area $6$, and shadows average a quarter of it.
\writeup (a) By Theorem~\ref{thm:cauchyspace}, the mean shadow is $\frac64=\frac32$. Directly, the shadow along $u$ has area $\frac12\sum_F|\langle n_F,u\rangle|=|u_1|+|u_2|+|u_3|$, whose mean is $3\cdot\frac12$. It is largest for
$u=\frac1{\sqrt3}(1,1,1)$, where it is $\sqrt3$; the shadow is then the regular hexagon formed by the projections of the six vertices other than the two on the diagonal. (b) Every shadow of the
sphere is a disc of area $\pi r^2$, and $\frac{4\pi r^2}4=\pi r^2$. A sphere is not a polytope, but it is the limit of inscribed polytopes, whose surface areas and shadows converge: the formula holds for
every convex body.
\end{solution}

\begin{solution}{ex:spherecrofton}
\idea Locally the sphere looks like the plane, and the computation of Crofton's formula goes through with the pole $u$ in place of $(\theta,p)$.
\writeup For a $C^1$ curve $c$ parametrized by arc length, $n(c\cap u^\perp)$ is the number of zeros of $s\mapsto\langle c(s),u\rangle$. By Lemma~\ref{lem:banach} in a form adapted to the sphere (the same proof), the measure of the poles $u$
whose great circle crosses a short arc from $c(s)$ to $c(s+\dd s)$ is that of a band of angular width $|\langle u,c'(s)\rangle|\dd s$ around the great circle $c(s)^\perp$, on which $u$ ranges over a great circle making the angle
$\phi$ with the tangent: $\langle u,c'(s)\rangle=\cos\phi$. The band has area $\int_0^{2\pi}|\cos\phi|\dd\phi\,\dd s=4\,\dd s$ out of $4\pi$, so $\int n\dd\sigma=\frac{4L}{4\pi}=\frac L\pi$. A great circle has $L=2\pi$, and almost every other great circle meets it
twice: $\frac{2\pi}\pi=2$.
\end{solution}

\section*{Sources and further reading}
\addcontentsline{toc}{section}{Sources and further reading}

J.~Bertrand, \emph{Calcul des probabilités} (Gauthier-Villars, 1889), pp.~4--5, where the chord problem and the verdict ``aucune des trois n'est fausse, aucune n'est exacte'' appear;
H.~Poincaré, \emph{Calcul des probabilités} (Carré, 1896; 2nd ed. Gauthier-Villars, 1912), the chapter on geometric probability, which introduces the invariant measures of lines and of motions and proves the formula of
Theorem~\ref{thm:poincare}; E.~T.~Jaynes, ``The well-posed problem'', \emph{Found.\ Phys.} 3 (1973), 477--492, with the straw-tossing experiment.

G.-L.~Leclerc de Buffon, \emph{Essai d'arithmétique morale}, in \emph{Histoire naturelle}, Supplément, vol.~4 (1777); P.-S.~Laplace, \emph{Théorie analytique des probabilités} (1812), for the grid; M.~W.~Crofton,
``On the theory of local probability'', \emph{Phil.\ Trans.\ R.\ Soc.\ London} 158 (1868), 181--199; J.-É.~Barbier, ``Note sur le problème de l'aiguille et le jeu du joint couvert'', \emph{J.\ Math.\ Pures Appl.} (2) 5 (1860),
273--286; A.~Cauchy, ``Note sur divers théorèmes relatifs à la rectification des courbes et à la quadrature des surfaces'', \emph{C.\ R.\ Acad.\ Sci.\ Paris} 13 (1841); H.~Steinhaus, ``Length,
shape and area'', \emph{Colloq.\ Math.} 3 (1954), 1--13; H.~Hadwiger, \emph{Vorlesungen über Inhalt, Oberfläche und Isoperimetrie} (Springer, 1957).

For more: L.~A.~Santaló, \emph{Integral Geometry and Geometric Probability} (Addison-Wesley, 1976; Cambridge, 2004), Chapters~1--7; D.~A.~Klain and G.-C.~Rota, \emph{Introduction to Geometric
Probability} (Cambridge University Press, 1997), which proves Hadwiger's theorem; and R.~Schneider and W.~Weil, \emph{Stochastic and Integral Geometry} (Springer, 2008).
